\documentclass[main.tex]{subfiles}
\begin{document}

\section{Exercise 4 – Bandgap Voltage Reference}

\subsection*{Operating Principle}

The bandgap reference generates a temperature-independent voltage by summing two complementary temperature-dependent voltages:

\textbf{CTAT (Complementary To Absolute Temperature):} The base-emitter voltage $V_{BE}$ decreases with temperature at $\approx -2$ mV/K.

\textbf{PTAT (Proportional To Absolute Temperature):} The differential voltage $\Delta V_{BE} = V_T \ln(M) = \frac{kT}{q}\ln(M)$ increases linearly with temperature, where $M$ is the emitter area ratio.

The reference voltage is:
$$V_{REF} = V_{BE} + G \cdot \Delta V_{BE}$$

The gain $G = R_2/R_1$ is chosen such that $\frac{\partial V_{REF}}{\partial T} = 0$, yielding $V_{REF} \approx 1.25$ V (silicon bandgap energy at 0K).

\subsection*{Circuit Implementation}

\textbf{Topology:} An op-amp forces equal voltages at two nodes, impressing $\Delta V_{BE}$ across a resistor to create PTAT current. This current flows through another resistor to generate the PTAT voltage component, which is summed with $V_{BE}$.

\textbf{Devices:} CMOS processes use parasitic PNP transistors (p+ emitter, n-well base, p-substrate collector).

\textbf{Startup:} Required to avoid the degenerate zero-current state. A small resistor injects initial current to establish proper biasing.

\subsection*{Advantages}

\begin{itemize}
    \item \textbf{Low TC:} Well-designed circuits achieve 10-20 ppm/°C
    \item \textbf{Supply independence:} High PSRR, insensitive to $V_{DD}$ variations
    \item \textbf{Process independence:} Based on fundamental physics (bandgap energy)
    \item \textbf{Fully integrated:} No external components needed
\end{itemize}

\subsection*{Disadvantages}

\begin{itemize}
    \item \textbf{Minimum supply:} Requires $V_{DD} \geq 1.5$ V (standard topology) due to headroom for $V_{REF} + V_{DS,sat}$. Low-voltage versions operate at $\approx 1$ V but with added complexity
    \item \textbf{Non-linearity:} Second-order temperature effects cause curvature over wide temperature ranges
    \item \textbf{Startup complexity:} Requires additional circuitry to ensure correct bias point
    \item \textbf{Trimming needed:} Process variations (±20-30\%) in absolute component values necessitate calibration for high accuracy
    \item \textbf{Noise:} Thermal noise from resistors and op-amp input noise require filtering
\end{itemize}

\subsection*{Applications \& Design Trade-offs}

\textbf{Ideal for:} ADC/DAC references, LDO regulators, temperature sensors requiring medium-voltage precision applications.

\textbf{Limitations:} Not suitable for low-voltage (<1.2V) applications without specialized topologies. High-precision applications require trimming, increasing production cost.

\textbf{Key design choice:} The resistor ratio $G$ and area ratio $M$ determine TC cancellation. Typical $M = 8-10$ provides adequate $\Delta V_{BE}$ while maintaining good matching.

\end{document}
